\documentclass{ximera}
\title{Looking back}
\begin{document}
\begin{abstract}
A short review of everything the course has covered.
\end{abstract}
\maketitle

That is the whole course. Before you put the notes away, it is worth looking back at how much ground you have covered:

\begin{itemize}
  \item \textbf{Functions} (Week 1): domain, range, the standard families of functions, and inverses --- including logarithms and the inverse trigonometric functions.
  \item \textbf{Limits} (Weeks 2--4): limits informally and precisely, one-sided and infinite limits, the limit laws, the squeeze theorem, continuity, and limits at infinity.
  \item \textbf{Derivatives} (Weeks 4--6): the derivative as a limit and as a rate of change, higher derivatives, and the rules --- power, exponential, product, quotient, chain --- together with implicit and logarithmic differentiation.
  \item \textbf{Applications of the derivative} (Weeks 8--10): related rates, exponential growth and decay, linear approximation and differentials, extrema and the mean value theorem, curve sketching, l'H\^opital's rule, and optimization.
  \item \textbf{Integration} (Weeks 11--13): antiderivatives, summation and Riemann sums, the definite integral and its properties, the fundamental theorem of calculus, indefinite integrals, and substitution.
\end{itemize}

If any of those is hazy, the fastest way to find out is to return to that section and redo the answer boxes with the notes closed.

\end{document}
